\section{Clifford-Error Analysis}
\label{sec:clifford_error_analysis}

In this section,
we consider in a general context how physical Clifford errors
affect logical states.
Specifically,
we consider afflicting a physical Clifford error on an encoded state,
then noiselessly measuring the stabiliser generators
(which we assume to be Paulis)
of the code.
Since the error in general is not Pauli,
these noiseless measurements in general are nondeterministic.
We may get a trivial syndrome ($\v s =\v 0$),
in which case the afflicted state is projected back into the codespace.
We are interested in two quantities:
(1) the probability to observe this trivial syndrome,
and (2) the fidelity of the resulting state given a trivial syndrome.

Conceptually,
our analysis uses that since Cliffords preserve commutation,
the Clifford error will transform the code
into a slightly different stabiliser code.
Calculating (1) involves counting
how many new stabilisers overlap with the old ones;
the time complexity of this is
polynomial in physical qubit count.
As for (2),
the result depends on the overlap between
the original projector defining the logical state
and the transformed projector.
\Cref{sec:acceptance_probability} answers (1) generally,
and \cref{sec:resulting_fidelity} answers (2)
in the case of T-state encoded in the colour code.

\subsection{Acceptance Probability}
\label{sec:acceptance_probability}
Here we present a general expression for
the probability $\pr(\v s =\v 0)$ to observe a trivial syndrome.
\begin{theorem}
\label{theorem:probability_trivial_syndrome}
Consider an arbitrary $k$-qubit mixed state $\rho$
encoded in an $\llbracket n, k\rrbracket$ Pauli stabiliser code,
which then suffers a physical Clifford error $E$.
Let $\mathcal G$ be a set of $n-k$ independent stabiliser
generators, which we measure noiselessly after the error.
The probability of a trivial syndrome is
\begin{equation}
\label{eq:probability_trivial_syndrome}
\pr(\v s =\v 0)
=\begin{cases*}
0, & if $`<E^\dagger \mathcal G E>
\cap \big(-`<\mathcal G>\big)
\neq \varnothing$, \\
2^{-s} \sum_{L} \omega_L \alpha_L,
& else,
\end{cases*}
\end{equation}
where the coefficients $\{\alpha_L\}_L$ define $\rho$ by
\cref{eq:pauli_decomposition},
and $s :=\rk`< \mathcal G\cup E^\dagger\mathcal G E > -(n -k)$ is
the number of new constraints imposed by the transformed stabilisers.
The sum is over all $L\in\mathcal P_k$ whose logical representative satisfies
\begin{equation}
\label{eq:logical_membership}
\overline L
\in \omega_L
`<E^\dagger \mathcal G E>
\langle\mathcal G\rangle
\end{equation}
for some sign $\omega_L\in\{\pm1\}$;
in the nonzero case, this sign is unique.
\end{theorem}

% \blue{
% \begin{theorem}
% Consider an arbitrary $k$-qubit mixed state $\rho$
% encoded in an $\llbracket n,k\rrbracket$ Pauli stabiliser code,
% which then suffers a physical Clifford error $E$.
% Let $\mathcal G$ be a set of $n-k$ independent stabiliser
% generators, which we measure noiselessly after the error.
% The probability of a trivial syndrome is
% \begin{equation}
%     \label{eq:probability_trivial_syndrome_zero}
%     \pr(\v s=\v 0)=0,
%     \qquad \text{if}\quad
%     \langle E^\dagger\mathcal G E\rangle
%     \cap \bigl(-\langle\mathcal G\rangle\bigr)
%     \neq \varnothing.
% \end{equation}

% Otherwise, let $\mathcal L_E\subseteq\mathcal P_k$ contain
% precisely those logical Paulis $L$ for which there exist
% a physical Pauli representative $\overline L$ of $L$
% and a sign $\omega_L\in\{+1,-1\}$ satisfying
% \begin{equation}
%     \label{eq:logical_membership}
%     \omega_L\overline L
%     \in \langle E^\dagger\mathcal G E\rangle.
% \end{equation}
% Then
% \begin{equation}
%     \label{eq:probability_trivial_syndrome}
%     \pr(\v s=\v 0)
%     =
%     2^{-s}
%     \sum_{L\in\mathcal L_E}\omega_L\alpha_L,
% \end{equation}
% where $\{\alpha_L\}_L$ are the Pauli coefficients of $\rho$
% defined in \cref{eq:pauli_decomposition}, and
% \[
%     s :=
%     \rk\langle\mathcal G\cup E^\dagger\mathcal G E\rangle
%     -(n-k)
% \]
% is the number of additional independent Pauli generators
% contributed by the transformed checks beyond the original
% stabiliser group.
% \end{theorem}
% Note that for each $L$ satisfying \cref{eq:logical_membership}, the sign $\omega_L$ is unique.
% }

\noindent
In words:
if the transformed stabiliser group contains
the negative of an original stabiliser,
the probability vanishes.
Otherwise,
a given logical Pauli component contributes to the probability if
any representative is a transformed stabiliser up to a sign.
We can interpret $2^{-s}$ as the acceptance probability for a maximally
mixed logical input, while the sum accounts for dependence on the
logical state.
We prove this statement in \cref{sec:clifford_error_analysis_proofs}
and present an $\mathcal O(n^3)$-time algorithm
to evaluate \cref{eq:probability_trivial_syndrome}
in \cref{sec:acceptance_probability_algorithm}.
In most cases considered in this paper,
the only logical that satisfies \cref{eq:logical_membership}
is identity.
Other logicals can contribute however,
if e.g.\ $E$ maps a stabiliser directly to that logical.

In the special case $E$ is a Pauli then either:
	$E$ flips at least one stabiliser in $\mathcal G$,
	in which case $\pr(\v s =\v 0) =0$,
	or $E^\dag \mathcal G E =\mathcal G$,
	in which case $s =0$ and only $L =\I$ satisfies
	\cref{eq:logical_membership} with $\omega_\I =1$,
	so $\pr(\v s =\v 0) =1$.

\begin{eg*}
To recap,
we have a T-state,
i.e.\ the state stabilised by $\H_+$,
encoded in the $\llbracket 7, 1, 3 \rrbracket$ hexagonal colour code,
which then suffers the Clifford error
$E =-(\H_+)_0 (\H_-)_{1 3}$.
$E$ transforms only $G_1$ and $G_2$ out of $\mathcal G$;
this can be seen from the support of $E$
and by using \cref{eq:h_pm_action,eq:d3_color_stabilizer_generators}.
Since $E=E^\dagger$,
\begin{equation}
(E^\dagger G_1E,\ldots,E^\dagger G_6E)=(
-\mathrm{YYXIXII},
\mathrm{IYXYIXI},
G_3,G_4,G_5,G_6).
\end{equation}
One checks that
$\langle E^\dagger \mathcal G E\rangle
\cap
\bigl(-\langle\mathcal G\rangle\bigr)=\varnothing$
and
$s=\rk\langle\mathcal G\cup E^\dagger\mathcal G E\rangle-n +k
=8-7+1=2$.
Also,
\begin{equation}
\proj{\T}
=\frac12\left(\I+\frac{\X+\Y}{\sqrt2}\right).
\end{equation}
Among the nonzero Pauli-basis coefficients,
only $L=\I$ satisfies \cref{eq:logical_membership},
with $\omega_\I=1$.
Therefore \cref{theorem:probability_trivial_syndrome}
gives $\pr(\v s=\v0) =2^{-s} \omega_\I \alpha_\I =1/4$.
\end{eg*}
\noindent
As an aside,
we recover the following result
that was implied by Aaronson and Gottesman~\cite[p~5]{Aaronson2004}.
\begin{cor}
The fidelity between two pure stabiliser states
whose Pauli stabiliser groups are
$\mathcal A$ and $\mathcal B$ respectively,
is $0$ if $\mathcal B \cap
(-\mathcal A)
\neq \varnothing$,
else $2^{-s}$
where $s := \rk\langle \mathcal A\cup \mathcal B \rangle -n$.
\end{cor}
\begin{proof}
Take $k =0$
in \cref{theorem:probability_trivial_syndrome},
and let $\mathcal A =`<\mathcal G>$
and $\mathcal B =`<E^\dag \mathcal G E>$
(i.e.\ the physical state stabilised by $\mathcal B$
is that of $\mathcal A$,
transformed by the Clifford $E^\dag$);
then the fidelity is $\pr(\v s =\v 0)$.
The density operator of the 0-qubit state is unity
i.e.\ $\alpha_1 =1$,
so \cref{eq:logical_membership}
reduces to $\overline{1} =\I \in \omega_1 \mathcal B \mathcal A$.
If $\omega_1 =-1$ then we are in the first branch of
\cref{eq:probability_trivial_syndrome} and the fidelity vanishes.
Else,
we are in the second branch and
$\sum_L \omega_L \alpha_L =\omega_1 \alpha_1 =1$.
\end{proof}

\subsection{Resulting Fidelity}
\label{sec:resulting_fidelity}
To answer the second question
(\textit{given a trivial syndrome,
what is the fidelity of the resulting state?}),
we consider the more specific scenario,
which we also prove in \cref{sec:clifford_error_analysis_proofs}:

\begin{prop}\label{prop:resulting_fidelity}
Consider a T-state encoded in an $\llbracket n, 1, d \rrbracket$
hexagonal colour code,
which then suffers a tensor product $E$ of Clifford errors from $`{\I, \H_\pm, \Z}$
up to global phase.
The operators $\H_+$, $\H_-$, and $\Z$
square to $\I$ and pairwise anticommute,
so $E$ either commutes or anticommutes with
$\overline{H_+}$ defined in \cref{eq:logical_H_is_transversal}.
If we noiselessly measure a minimal set $\mathcal G$ of stabiliser generators
and observe a trivial syndrome,
the fidelity of the resulting state is
one (zero) if $E$ commutes (anticommutes) with
$\overline{H_+}$;
determining this takes $\mathcal O(n)$ time.
\end{prop}

\begin{eg*}
For $E =-(\H_+)_0(\H_-)_{1 3}$,
the only tensor factor that anticommutes with the corresponding
factor of $\overline{H_+}=(\H_+)_{0 2 3 6}(\H_-)_{1 4 5}$
is at qubit 3.
Thus $E$ anticommutes with $\overline{H_+}$, so the fidelity is zero.
\end{eg*}

Lastly,
we note there is an alternative analysis
(to \cref{theorem:probability_trivial_syndrome,prop:resulting_fidelity})
that treats the Clifford error as a superposition of Paulis
instead of focusing on how it transforms the stabilisers.
This recovers the same results
and is potentially more intuitive,
but is slower:
computing $\pr(\v s =\v 0)$ in $\mathcal O(2^n)$ time.
In \cref{sec:alternative_clifford_error_analysis}
we walk through \cref{eg:d3_color_code_acceptance_probability}
using this alternative analysis.